\DocumentMetadata{
  pdfversion=2.0,
  pdfstandard=ua-2,
  lang=en-US,
  tagging=on,
  tagging-setup={math/setup=mathml-SE,tabsorder=structure}
}
\documentclass[11pt,letterpaper]{article}
\usepackage[margin=0.68in,includefoot]{geometry}
\usepackage{newcomputermodern}
\usepackage{amsmath,amscd,mathtools}
\usepackage{tikz,adjustbox}
\providecommand{\square}{\mdlgwhtsquare}
\usepackage[hidelinks]{hyperref}
\hypersetup{
  pdftitle={Algebra First-Year Exam, January 2024, Part 2},
  pdfauthor={University of Florida Department of Mathematics},
  pdfsubject={Graduate mathematics examination},
  pdfdisplaydoctitle=true,
  bookmarksopen=true
}
\setcounter{secnumdepth}{0}
\setlength{\parindent}{0pt}
\setlength{\parskip}{0.28em}
\setlength{\emergencystretch}{3em}
\setlength{\leftmargini}{2.15em}
\setlength{\leftmarginii}{2.65em}
\setlength{\leftmarginiii}{3em}
\pagestyle{empty}
\raggedbottom
\allowdisplaybreaks
\NewTaggingSocketPlug{block/list/label}{examproblem}{%
  \tagstructbegin{tag=Lbl}\tagstructend%
  \tagstructbegin{tag=\LBody}%
  \tagstructbegin{tag=\ProblemHeadingTag,title-o={\ProblemHeadingText}}%
  \tagmcbegin{tag=\ProblemHeadingTag,actualtext={\ProblemHeadingText}}%
  #2%
  \tagmcend%
  \tagstructend%
}
\newenvironment{examproblems}[2]{%
  \def\ProblemHeadingTag{#2}%
  \AssignTaggingSocketPlug{block/list/label}{examproblem}%
  \begin{enumerate}%
  \setcounter{enumi}{#1}%
  \renewcommand{\labelenumi}{\textbf{\theenumi.}}%
  \setlength{\topsep}{0.35em}%
  \setlength{\partopsep}{0pt}%
  \setlength{\itemsep}{0.48em}%
  \setlength{\parsep}{0.18em}%
}{\end{enumerate}}
\newenvironment{examsubparts}{%
  \AssignTaggingSocketPlug{block/list/label}{default}%
  \begin{enumerate}%
  \setlength{\topsep}{0.18em}%
  \setlength{\partopsep}{0pt}%
  \setlength{\itemsep}{0.16em}%
  \setlength{\parsep}{0.1em}%
}{\end{enumerate}}
\newenvironment{examinstructions}{%
  \par\begingroup\itshape
}{%
  \par\endgroup\vspace{0.18em}
}
\makeatletter
\renewcommand\subsection{\@startsection{subsection}{2}{0pt}%
  {-1.25ex plus -.25ex minus -.1ex}%
  {0.55ex plus .1ex}%
  {\normalfont\large\bfseries}}
\renewcommand{\@maketitle}{%
  \newpage\null\vskip -1em%
  {\footnotesize\bfseries
    \href{https://www.ufl.edu/}{University of Florida}, \href{https://math.ufl.edu/}{Department of Mathematics}\par}%
  \vskip -0.5em%
  \tagstructbegin{tag=H1,title={Algebra First-Year Exam, January 2024, Part 2}}%
  \tagpdfparaOff%
  \tagmcbegin{tag=H1}%
    {\LARGE\bfseries\href{https://gma.math.ufl.edu/exams/algebra-first-year/}{Algebra First-Year Exam}, January 2024, Part 2\par}%
  \tagmcend%
  \tagpdfparaOn%
  \tagstructend%
  \vskip 0.25em%
  \tagmcbegin{artifact}%
    \hrule height 1pt%
  \tagmcend%
  \vskip 1em%
}
\makeatother
\title{Algebra First-Year Exam, January 2024, Part 2}
\author{}
\date{}
\begin{document}
\maketitle
\thispagestyle{empty}
\begin{examinstructions}
Answer four problems. Write your answers clearly in complete English sentences. You may quote results (within reason) as long as you state them clearly.
\end{examinstructions}
\begin{examproblems}{0}{H2}
\def\ProblemHeadingText{Problem 1.}
\item
Suppose~\(k\) is a field.
\begin{examsubparts}
\item[\textnormal{(i)}]
If \(A=k[X,Y]/(X^2Y^3)\), find all minimal prime ideals of~\(A\).
\item[\textnormal{(ii)}]
Show that \(B=k[X,Y]/(X^2,Y^3)\) has a unique prime ideal, and that this prime ideal is maximal. What is it?
\end{examsubparts}
\def\ProblemHeadingText{Problem 2.}
\item
Which of the following polynomials are irreducible? Explain.
\begin{examsubparts}
\item[\textnormal{(a)}]
\(X^3+X+1\) in \(\mathbb{F}_2[X]\);
\item[\textnormal{(b)}]
\(X^3+X+1\) in \(\mathbb{Q}[X]\);
\item[\textnormal{(c)}]
\(X^5+10X^2+25X+5\) in \(\mathbb{Q}[X]\);
\end{examsubparts}
\def\ProblemHeadingText{Problem 3.}
\item
Suppose~\(R\) is a ring with identity and~\(M\) is a nonzero left~\(R\)-module. Let \(\operatorname{End}_R(M)\) be the ring of~\(R\)-linear maps \(M\to M\). Show that if~\(0\) and~\(M\) are the only~\(R\)-submodules of~\(M\), then \(\operatorname{End}_R(M)\) is a division ring.
\def\ProblemHeadingText{Problem 4.}
\item
Let~\(R\) be an integral domain.
\begin{examsubparts}
\item[\textnormal{(i)}]
Define what it means for an element of~\(R\) to be irreducible, and to be prime.
\item[\textnormal{(ii)}]
Show that a prime element is irreducible.
\end{examsubparts}
\def\ProblemHeadingText{Problem 5.}
\item
Suppose~\(F\) is a field, \(V\) and~\(W\) are~\(F\)-vector spaces of finite dimension and \(f:V\to W\) is a surjective linear map, and let \(U=\operatorname{Ker}(f)\). Show using the definitions that \(\dim_F(V)=\dim_F(W)+\dim_F(U)\).
\def\ProblemHeadingText{Problem 6.}
\item
Suppose \(L/K\) is an extension of fields. Show that \(\alpha\in L\) is algebraic over~\(K\) if and only if there is a nonzero sub-\(K\)-vector space \(V\subseteq L\) of finite dimension such that \(\alpha V\subseteq V\) (for “if,” pick a nonzero \(v\in V\) and consider \(\alpha^n v\) for \(n\in\mathbb{N}\)).
\def\ProblemHeadingText{Problem 7.}
\item
Find representatives of each conjugacy class of \(GL_4(\mathbb{F}_2)\) of order~\(3\) or~\(5\) (the order of a conjugacy class is the order of any element of it).
\end{examproblems}
\end{document}
